\documentclass{article}

\usepackage[utf8]{inputenc}
\usepackage[T1]{fontenc}
\usepackage[spanish]{babel}
\usepackage{amsmath}
\usepackage{amssymb}
\usepackage{geometry}
\usepackage{tikz}
\usepackage{booktabs}
\usepackage{xcolor}
\usepackage{array}
\geometry{margin=2.5cm}

\usetikzlibrary{automata, positioning, arrows.meta, fit, backgrounds}

\tikzset{
  every state/.style={
    draw, rounded corners=6pt,
    minimum width=2.2cm, minimum height=0.75cm,
    font=\small, align=center,
    fill=gray!8
  },
  initial text={},
  every edge/.style={draw, -{Stealth[length=5pt]}, font=\scriptsize},
  loop above/.append style={looseness=5},
  loop below/.append style={looseness=5},
}

\title{Modelos DEVS del sistema de bomba de infusión}
\date{}

% ── helpers ────────────────────────────────────────────────────────────────────
\newcommand{\dint}{\delta_{\mathrm{int}}}
\newcommand{\dext}{\delta_{\mathrm{ext}}}
\newcommand{\ta}{\mathit{ta}}

\begin{document}

\maketitle

\tableofcontents
\newpage

% ==============================================================================
\section{Notación común}

Cada modelo atómico DEVS se especifica mediante la tupla
\[
  M = \langle X,\; Y,\; S,\; \dint,\; \dext,\; \lambda,\; \ta \rangle.
\]
Se adoptan las siguientes convenciones a lo largo del documento:

\begin{itemize}
  \item Los puertos de entrada y salida se distinguen con un índice entero
        no negativo: $(v, p)$ donde $v$ es el valor y $p \in \mathbb{N}_0$
        es el número de puerto.
  \item En las transiciones externas, $e$ denota el tiempo transcurrido desde
        la última transición.
  \item $\mathbb{R}_0^+$ denota los reales no negativos incluyendo $\infty$.
  \item Los estados pasivos utilizan $\sigma = \infty$ como avance temporal.
\end{itemize}

% ==============================================================================
\section{Generador de órdenes médicas}

El generador produce órdenes médicas periódicas cuyo caudal e intervalo
de tiempo se muestrean de las distribuciones definidas en el documento
de modelado estadístico, parametrizadas por el grado de criticidad $k$.

\subsection{Definición formal}

\[
  M_{\mathrm{gen}} = \langle X,\; Y,\; S,\; \dint,\; \dext,\; \lambda,\; \ta \rangle
\]

\begin{align*}
  X &= \varnothing \\[4pt]
  Y &= \bigl\{\,(c,\,t,\;0)\;\big|\;
        c\in\{0,50,100,150,200\},\;
        t\in\{2,4,6,8,12\}\bigr\} \\[4pt]
  S &= \bigl\{\,(k,\,c,\,t)\;\big|\;
        k\in[0,1],\;
        c\in\{0,50,100,150,200\},\;
        t\in\{2,4,6,8,12\}\bigr\}
\end{align*}

\paragraph{Transición interna.}
Genera una nueva orden a partir del grado de criticidad actual:
\[
  \dint(k,\,c,\,t)
  \;=\;
  \bigl(\,k,\;\mathrm{Caudal}(k),\;\mathrm{Tiempo}(k)\,\bigr).
\]

\paragraph{Transición externa.}
El modelo no recibe eventos externos:
\[
  \dext = \varnothing.
\]

\paragraph{Función de salida.}
Emite el caudal y el tiempo de la orden vigente:
\[
  \lambda(k,\,c,\,t) \;=\; (c,\,t,\;0).
\]

\paragraph{Avance temporal.}
\[
  \ta(k,\,c,\,t) \;=\; t.
\]

\paragraph{Estado inicial.}
\[
  s_0 \;=\; \bigl(\,k_0,\;\mathrm{Caudal}(k_0),\;\mathrm{Tiempo}(k_0)\,\bigr),
\]
donde $k_0$ es el grado inicial de criticidad del paciente.

\subsection{Diagrama de estados}

\begin{center}
\begin{tikzpicture}[node distance=3.5cm]
  \node[state, initial left] (gen) {$(k,\,c,\,t)$};
  \path (gen) edge[loop right, looseness=5]
    node[right] {$\dint\;/\;\lambda{:}\;(c,t,0)$} (gen);
\end{tikzpicture}
\end{center}

El modelo tiene un único estado paramétrico. Cada transición interna
actualiza $(c,t)$ mediante muestreo y emite la orden antes de actualizar.

% ==============================================================================
\section{Sensor de flujo}

El sensor mide el caudal real administrado por la bomba y emite una
lectura cada segundo.

\subsection{Definición formal}

\begin{align*}
  X &= \{\,(c,\;0)\;\mid\; c\in[0,200]\,\} \\[4pt]
  Y &= \{\,(c,\;0)\;\mid\; c\in[0,200]\,\} \\[4pt]
  S &= \{\,(c,\;\sigma)\;\mid\; c\in[0,200],\;\sigma\in\mathbb{R}_0^{+}\,\}
\end{align*}

El estado $s = (c,\,\sigma)$ contiene el último caudal recibido del
actuador y el tiempo restante $\sigma$ hasta la próxima medición.

\paragraph{Transición externa.}
Actualiza el caudal medido sin modificar el tiempo restante:
\[
  \dext\bigl((c,\,\sigma),\;e,\;(x,0)\bigr)
  \;=\; (x,\;\sigma - e).
\]

\paragraph{Transición interna.}
Emite la medición y programa la siguiente un segundo después:
\[
  \dint(c,\,\sigma) \;=\; (c,\;1).
\]

\paragraph{Función de salida.}
\[
  \lambda(c,\,\sigma) \;=\; (c,\;0).
\]

\paragraph{Avance temporal.}
\[
  \ta(c,\,\sigma) \;=\; \sigma.
\]

\paragraph{Estado inicial.}
\[
  s_0 \;=\; (0,\;\infty).
\]
El sensor permanece pasivo hasta recibir la primera lectura del actuador.

\subsection{Diagrama de estados}

\begin{center}
\begin{tikzpicture}[node distance=4cm]
  \node[state, initial left] (pasivo) {pasivo\\$(c,\infty)$};
  \node[state, right of=pasivo] (activo) {activo\\$(c,\sigma)$};

  \path (pasivo) edge[bend left=20]
      node[above] {\scriptsize $\dext{:}\;(x,0)\;\big/\;(x,\,\infty{-}0)$}
      (activo)
    (activo) edge[loop right, looseness=5]
      node[right] {\scriptsize $\dint\;/\;\lambda{:}\;(c,0)$}
      (activo)
    (activo) edge[loop below, looseness=5]
      node[below] {\scriptsize $\dext{:}\;(x,0)\;\big/\;(x,\,\sigma{-}e)$}
      (activo);
\end{tikzpicture}
\end{center}

% ==============================================================================
\section{Controlador}

El controlador recibe órdenes médicas, lecturas del sensor, señales de
fin de bolsa y confirmaciones del enfermero; emite ajustes a la bomba
y dispara alarmas.

\subsection{Definición formal}

\begin{align*}
  X &= \{\mathrm{ordenMedica}\}\!\times\!\{0\}
     \;\cup\; \{\mathrm{sensorFlujo}\}\!\times\!\{1\}
     \;\cup\; \{\mathrm{finBolsa}\}\!\times\!\{2\}
     \;\cup\; \{\mathrm{confirmacionEnfermero}\}\!\times\!\{3\}
  \\[4pt]
  Y &= \{\mathrm{alarmaBaja},\mathrm{alarmaMedia},\mathrm{alarmaCritica}\}
       \!\times\!\{0\}
     \;\cup\;
     \{\mathrm{ajustarCaudal},\mathrm{detenerBomba}\}
       \!\times\!\{1\}
  \\[4pt]
  S &= \underbrace{\mathbb{R}_{\ge 0}}_{\mathit{uCM}}
     \times
     \underbrace{\mathbb{R}_{\ge 0}}_{\mathit{cO}}
     \times
     \underbrace{\mathbb{N}_0}_{\mathit{tol}}
     \times
     \underbrace{E_{\mathrm{ctrl}}}_{\mathit{est}}
     \times
     \underbrace{\mathbb{R}_0^+ \cup \{\infty\}}_{\mu}
\end{align*}

donde

\[
E_{\mathrm{ctrl}}
=
\{
\mathrm{normal},
\mathrm{bolsaBaja},
\mathrm{esperandoDetencion},
\mathrm{estadoCritico}
\}
\]

y los componentes del estado representan:

\begin{itemize}
\item \(uCM\): último caudal medido.
\item \(cO\): caudal objetivo.
\item \(tol\): contador de segundos consecutivos fuera de tolerancia.
\item \(est\): estado lógico del controlador.
\item \(\mu\): tiempo hasta la próxima transición interna.
\end{itemize}

\paragraph{Estado inicial.}

\[
s_0=
(0,\;0,\;0,\;\mathrm{normal},\;\infty)
\]

\paragraph{Función auxiliar.}

\[
\mathrm{no\_tolerable}(cO,uCM)
=
\frac{|cO-uCM|}{cO}
>
0.10
\]

\paragraph{Transición externa.}

Recepción de una nueva orden médica:

\[
\dext
((uCM,cO,tol,est,\mu),
e,
((c,t),0))
=
(uCM,c,tol,est,0)
\]

Recepción de una medición del sensor:

\[
\dext
((uCM,cO,tol,est,\mu),
e,
(x,1))
=
\begin{cases}

(x,cO,tol+1,estadoCritico,0)
&
\text{si }
\mathrm{tol > k_{\mathrm{crit}}\times 5
}
\\[2ex]


(x,cO,tol+1,est,0)
&
\text{si }
\mathrm{no\_tolerable}(cO,x)
\\[2ex]
(x,cO,0,est,0)
&
\text{en otro caso}
\end{cases}
\]

Recepción de señal de fin de bolsa:

\[
\dext
((uCM,cO,tol,est,\mu),
e,
(\mathrm{finBolsa},2))
=
(uCM,cO,tol,\mathrm{bolsaBaja},0)
\]

Recepción de confirmación del enfermero:

\[
\dext
((uCM,cO,tol,est,\mu),
e,
(\mathrm{confirmacionEnfermero},3))
=
(uCM,cO,0,\mathrm{normal},\infty)
\]

\paragraph{Transición interna.}

Luego de emitir una alarma por bolsa baja:

\[
\dint
(uCM,cO,tol,\mathrm{bolsaBaja},0)
=
(uCM,cO,tol,\mathrm{esperandoDetencion},60)
\]

Luego de transcurridos los 60 segundos:

\[
\dint
(uCM,cO,tol,\mathrm{esperandoDetencion},60)
=
(uCM,cO,tol,\mathrm{normal},\infty)
\]

En cualquier otro caso:

\[
\dint
(uCM,cO,tol,\mathrm{normal},\mu)
=
(uCM,cO,tol,\mathrm{normal},\infty)
\]

\paragraph{Función de salida.}

Los casos se evalúan en el orden indicado.

\[
\lambda(uCM,cO,tol,est,\mu)
=
\begin{cases}

(\mathrm{detenerBomba},1)
&
\text{si }
estado = estadoCritico

\\[2ex]

(\mathrm{detenerBomba},1)
&
\text{si }
est=\mathrm{esperandoDetencion}
\\[2ex]



(\mathrm{detenerBomba},1)
&
\text{si }
cO=0
\\[2ex]

(\mathrm{alarmaCritica},0)
&
\text{si }
tol > k_{\mathrm{crit}}\times 5
\\[2ex]

(\mathrm{alarmaMedia},0)
&
\text{si }
tol > 5
\\[2ex]

(\mathrm{ajustarCaudal}(cO-uCM),1)
&
\text{si }
uCM \neq cO
\land
est=\mathrm{normal}
\\[2ex]

(\mathrm{ajustarCaudal}(0),1)
&
\text{si }
uCM = cO
\end{cases}
\]



\paragraph{Avance temporal.}

\[
ta(uCM,cO,tol,est,\mu)
=
\mu
\]

\section{Bomba de infusión}

La bomba modifica el caudal real en respuesta a comandos del controlador.

\subsection{Definición formal}

\begin{align*}
  X &= \{\mathrm{ajustarBomba},\;\mathrm{detenerBomba}\}\times\{1\} \\[4pt]
  Y &= \{\,(c,\;0)\;\mid\;c\in[0,200]\,\} \\[4pt]
  S &= \underbrace{[0,200]}_{\mathit{cR}}
     \times \underbrace{\{\mathrm{detenida},\,\mathrm{corriendo}\}}_{\mathit{est}}
     \times \underbrace{\mathbb{R}_0^+}_{\mu}
\end{align*}

\paragraph{Transición externa.}

\begin{align*}
\dext&\bigl((cR,\,est,\,\mu),\;e,\;(\mathrm{detenerBomba},\,1)\bigr)
  \;=\; (0,\;\mathrm{detenida},\;0)
\\[6pt]
\dext&\bigl((cR,\,est,\,\mu),\;e,\;(\mathrm{ajustarBomba}(\mathit{dif}),\,1)\bigr)
  \;=\; \Bigl(cR + \frac{\mathit{dif}}{\mathrm{aleatorio}()},\;
              \mathrm{corriendo},\;
              \mathrm{random\_truncado}(0.5)\Bigr)
\end{align*}

\noindent
donde $\mathrm{aleatorio}()$ modela la inercia del mecanismo físico de la
bomba ($\mathrm{aleatorio}() \in (0,1]$).

\paragraph{Transición interna.}
Una vez transcurrido el retardo $\mu$, la bomba publica su caudal real
y queda pasiva:
\[
  \dint(cR,\,est,\,\mu) \;=\; (cR,\;est,\;\infty).
\]

\paragraph{Función de salida.}
Emite el caudal real actualizado:
\[
  \lambda(cR,\,est,\,\mu) \;=\; (cR,\;0).
\]

\paragraph{Avance temporal.}
\[
  \ta(cR,\,est,\,\mu) \;=\; \mu.
\]

\paragraph{Estado inicial.}
\[
  s_0 \;=\; (0,\;\mathrm{detenida},\;\infty).
\]

\subsection{Diagrama de estados}

\begin{center}
\begin{tikzpicture}[node distance=5cm]
  \node[state, initial left] (det) {detenida\\$cR=0$};
  \node[state, right of=det]  (cor) {corriendo\\$cR>0$};

  \path
    (det) edge[bend left=25]
      node[above] {\scriptsize $\dext{:}$ ajustar}
      (cor)
    (cor) edge[bend left=25]
      node[below] {\scriptsize $\dext{:}$ detener $\Rightarrow (0,\text{det},0)$}
      (det)
    (cor) edge[loop right, looseness=5]
      node[right] {\scriptsize $\dint\;/\;\lambda{:}\;(cR,0)$}
      (cor)
    (cor) edge[loop above, looseness=5]
      node[above] {\scriptsize $\dext{:}$ ajustar}
      (cor);
\end{tikzpicture}
\end{center}

% ==============================================================================
\section{Módulo de alarmas}

El módulo recibe señales de alarma del controlador, las emite hacia el
entorno (interfaz de usuario / personal médico) y gestiona el ciclo de
espera de confirmación para las alarmas críticas.

\subsection{Definición formal}

\begin{align*}
  X &= \{\mathrm{alarmaBaja},\,\mathrm{alarmaMedia},\,\mathrm{alarmaCritica}\}
       \!\times\!\{0\}
     \;\cup\; \{\mathrm{confirmacionEnfermero}\}\!\times\!\{1\} \\[4pt]
  Y &= \{\mathrm{alarmaBaja},\,\mathrm{alarmaMedia},\,\mathrm{alarmaCritica}\}
       \!\times\!\{0\} \\[4pt]
  S &= E_{\mathrm{alarmas}} \times \mathbb{R}_0^+
\end{align*}

El conjunto de estados discretos es:
\[
  E_{\mathrm{alarmas}} = \{
    \mathrm{noAlarma},\;
    \mathrm{alarmaBaja},\;
    \mathrm{alarmaMedia},\;
    \mathrm{alarmaCritica},\;
    \mathrm{espLargo},\;
    \mathrm{espCorto},\;
    \mathrm{repetirCritica}
  \}
\]

El estado general se expresa como $s = (\mathit{est},\,\sigma)$ donde
$\sigma$ es el tiempo hasta la próxima transición interna.

\paragraph{Estado inicial.}
\[
  s_0 \;=\; (\mathrm{noAlarma},\;\infty).
\]

\paragraph{Transición externa.}
\[
\dext\bigl((\mathit{est},\,\sigma),\;e,\;(x,\,p)\bigr) =
\begin{cases}
  (\mathrm{alarmaCritica},\;0)
    & p=0,\;x=\mathrm{alarmaCritica} \\[4pt]
  (x,\;0)
    & p=0,\;x\in\{\mathrm{alarmaBaja},\,\mathrm{alarmaMedia}\} \\[4pt]
  (\mathrm{noAlarma},\;\infty)
    & p=1,\;\mathit{est}\in\{\mathrm{espLargo},\,\mathrm{espCorto},\,\mathrm{repetirCritica}\} \\[4pt]
  (\mathit{est},\;\sigma - e)
    & \text{en otro caso}
\end{cases}
\]

\paragraph{Transición interna.}
\[
\dint(\mathit{est},\,\sigma) =
\begin{cases}
  (\mathrm{noAlarma},\;\infty)       & \mathit{est}\in\{\mathrm{alarmaMedia},\,\mathrm{alarmaBaja}\} \\[4pt]
  (\mathrm{espLargo},\;30)           & \mathit{est} = \mathrm{alarmaCritica} \\[4pt]
  (\mathrm{repetirCritica},\;0)      & \mathit{est} = \mathrm{espLargo} \\[4pt]
  (\mathrm{espCorto},\;10)           & \mathit{est} = \mathrm{repetirCritica} \\[4pt]
  (\mathrm{repetirCritica},\;0)      & \mathit{est} = \mathrm{espCorto}
\end{cases}
\]

\paragraph{Función de salida.}
\[
\lambda(\mathit{est},\,\sigma) =
\begin{cases}
  (\mathrm{alarmaCritica},\;0)
    & \mathit{est}\in\{\mathrm{alarmaCritica},\,\mathrm{repetirCritica}\} \\[4pt]
  (\mathit{est},\;0)
    & \mathit{est}\in\{\mathrm{alarmaMedia},\,\mathrm{alarmaBaja}\}
\end{cases}
\]

\paragraph{Avance temporal.}
\[
  \ta(\mathit{est},\,\sigma) \;=\; \sigma.
\]

\subsection{Diagrama de estados}

\begin{center}
\begin{tikzpicture}[
    node distance=2.2cm and 3.2cm,
    every state/.append style={minimum width=2.4cm}
  ]

  \node[state, initial above] (noA)  {noAlarma};
  \node[state, below left=of noA]  (aB)   {alarmaBaja};
  \node[state, below=of noA]       (aM)   {alarmaMedia};
  \node[state, below right=of noA] (aC)   {alarmaCritica};

  \node[state, below=2cm of aC]    (eL)   {espLargo\\$\sigma{=}30$};
  \node[state, below=2cm of eL]    (rC)   {repetirCritica};
  \node[state, right=of rC]        (eC)   {espCorto\\$\sigma{=}10$};

  % entradas desde noAlarma
  \path (noA) edge[bend right=10] node[left]  {\scriptsize alarmaBaja}   (aB)
        (noA) edge               node[right] {\scriptsize alarmaMedia}  (aM)
        (noA) edge[bend left=10] node[right] {\scriptsize alarmaCritica}(aC);

  % expiración de baja/media → noAlarma
  \path (aB)  edge[bend right=20] node[left]  {\scriptsize $\dint$} (noA)
        (aM)  edge               node[left]  {\scriptsize $\dint$} (noA);

  % ciclo crítico
  \path (aC)  edge node[right] {\scriptsize $\dint$: espLargo(30)} (eL)
        (eL)  edge node[right] {\scriptsize $\dint$: repetir}      (rC)
        (rC)  edge node[above] {\scriptsize $\dint$: espCorto(10)} (eC)
        (eC)  edge[bend left=30]
              node[right] {\scriptsize $\dint$: repetir}           (rC);

  % confirmación → noAlarma
  \path (eL)  edge[bend left=55, dashed]
              node[above left] {\scriptsize confirm.} (noA)
        (eC)  edge[bend right=60, dashed]
              node[below right] {\scriptsize confirm.} (noA)
        (rC)  edge[bend right=40, dashed]
              node[below] {\scriptsize confirm.} (noA);

\end{tikzpicture}
\end{center}

Las transiciones punteadas representan $\dext$ con puerto $p=1$
(confirmación del enfermero), que reinicia el módulo a \textit{noAlarma}.

\end{document}